\begin{frame}{역사: ``경험을 다시 쓴다''는 아이디어의 뿌리}
  \begin{itemize}
    \item 경험 재현은 DQN이 처음 발명한 것이 아니다 ---
      ``과거 경험을 다시 쓴다''는 정신은 1990년대 표 기반 문헌에서
      이미 여러 갈래:
    \item \kb{Dyna}(Sutton, 1991, Week 7.3): 실제 경험과
      \textbf{모델 시뮬레이션} 경험을 같은 스트림에 섞어 TD 갱신
      (``여러 출처의 경험을 하나로 합쳐 재사용''의 원형).
    \item \kb{Q($\lambda$)}(Week 7.2): 적격 이력으로 한 경험을
      \textbf{여러 스텝에 걸쳐} 재사용.
    \item \kb{배치(오프라인) RL}: 모은 데이터를 \textbf{여러 에포크} 반복 사용.
  \end{itemize}
  \vskip0.4em
  \begin{block}{DQN = 규모의 승리}
    이 뿌리 위에 \textbf{신경망 $+$ 대규모 버퍼($10^4 \sim 10^6$) $+$
    버퍼에서 무작위 재샘플링}을 태워, Atari의 \textbf{수백만 프레임}에서
    상관 문제를 ``단순한 무작위 추출 한 줄''로 해결.
  \end{block}
  {\footnotesize 원논문 DQN: 버퍼 $10^6$ 프레임, 배치 32,
  \textbf{4스텝마다 1회} 학습 --- ``4스텝마다''는 매 프레임 학습보다
  상관(4프레임은 거의 같은 화면)을 \textbf{더 강하게} 끊으려는 추가
  상관 제거.}
\end{frame}
